\chapter{METHODOLOGY}
\section{Research Design}
This thesis uses a qualitative, observation-informed technical assessment. Its first information stream is a practitioner account supplied through the project lead and subsequently expressed in engineering English. Its second stream is the existing peer-reviewed literature on rework, and its third is relevant, publicly identifiable technical guidance for waterproofing, concrete pavements and masonry. The streams serve different functions: the practitioner material identifies concerns; technical sources identify mechanisms worth testing; documentary project records, where obtained, establish whether a specific corrective rework event occurred.

The account was described as arising from approximately five years of practice in Pokhara, comprising three years of site engineering and two years of project management. That description is not attributed to the student because the observer's identity, affiliation and authority to contribute research evidence have not been confirmed. The original mixed Nepali--English statement and edited technical translation have been retained within the governed research record. Potential ambiguity in the original negative wording on wall toothing, bands and brick wetting requires confirmation by the speaker. The English technical interpretation is used here as a research-input narrative, not as a verified project-level transcript or structured interview. No interviews or questionnaires were conducted or assumed in this revision.

\section{Units and Research Boundary}
Three levels of information are distinguished. A \emph{reported construction concern} is a statement about conditions encountered in practice; a \emph{suspected rework case} links a traceable site and previous work to a possible correction; and a \emph{verified rework event} requires evidence of previously executed work together with a subsequent corrective intervention. No number of sites, visits, interviewees or admitted cases is assumed. An observation may concern workmanship quality even when there is no evidence that rework followed.

The project scope remains RCC building construction in Pokhara Metropolitan City. W01 waterproofing and W03 plaster/masonry are relevant associated building activities. W02 concrete pavement is included in case analysis only if the pavement is part of a documented building-development contract, site access, courtyard or similar project component. Stand-alone public road work remains outside the admitted case population unless the research scope is formally broadened. Discussion of pavement engineering mechanisms does not by itself admit a municipal road as a thesis case.

\section{Research-Input Mapping}
The supplied account was mapped into three categories without claiming that each remark represents a distinct project occurrence. The statements cover: missing or deficient wet-area waterproofing, terrace falls, joints, temporary support openings and vertical membrane details (W01); cracks in reported 4--6 inch PCC pavement, curing and questions about reinforcement (W02); and perpendicular brick junctions, opening bands, brick moisture preparation, masonry curing and associated plaster cracking (W03).

\begin{table}[htbp]
\caption{Source-to-analysis sequence for the reported material}
\label{tab:sourceflow}
\centering
\begin{tabularx}{\textwidth}{@{}L{3.1cm}X@{}}
\toprule Stage & Procedure \\ \midrule
Account transcription & Preserve the source language and any verified translation, identify the observer and date when available.\\
Topic mapping & Assign the technical subject W01, W02 or W03; a statement is not counted as a distinct case.\\
Mechanism screening & List credible alternative mechanisms and the test, drawing or site record needed to distinguish them.\\
Technical comparison & Check authentic project specifications and applicable codes; use international guidance only as technical context.\\
Event classification & Admit rework only after original executed work and documented corrective activity are both established.\\
Consequence assessment & Assess attributable quantities, rates and delay only from supporting records; otherwise leave unquantified.\\ \bottomrule
\end{tabularx}
\end{table}

\section{Technical Reference Selection}
The ten pre-existing peer-reviewed rework papers are retained. Supplementary technical publications were selected from recognized institutional or manufacturer sources where the established corpus does not address a mechanism. Examples include the DUDBC building-code publication catalogue, the Nepal Department of Roads' 2082 amendment, FHWA pavement guidance, the Tile Council of North America standards catalogue, waterproofing-system technical information and brick absorption research \parencite{dudbc2024,dor2082,fhwa2013,tcna2023,bailey1990}. Institutional or manufacturer guidance is explicitly distinguished from local empirical evidence; no unverified code clause or minimum dimension is stated as universally mandatory.

\section{Operational Definitions and Verification}
A defect is a condition suspected or demonstrated to deviate from a requirement. A mechanism is the physical process by which a condition may arise or progress. A cause is a verified contributing event or decision, not merely the visible symptom. Rework is an actual corrective intervention affecting work already carried out. Maintenance, planned finishing, first installation and routine curing are not counted automatically as rework.

A potential field event should be assigned a unique, anonymized reference. The observation record should include the element and location, date of reported condition, original approved requirement, evidence that the original work was carried out, a subsequent instruction or action, source records and photographs where consent exists, and the verification decision. Where the project identity, original-work evidence or corrective action cannot be matched, the case should remain \emph{unverified}. An event can be verified while its cause remains undetermined and its cost cannot be calculated.

\section{Analysis Procedure}
The first stage is descriptive: statements from the practitioner account are organized under the three subjects without asserting a frequency. The second stage is technical: each concern is compared with alternative mechanisms and relevant published guidance. The third stage is comparative: the recurrent themes of supervision, detailing, workmanship preparation, curing and verification are examined across subjects. The fourth stage addresses implications: possible corrective interventions are identified, clearly labelled as potential unless an actual repair record is available. Findings in Chapter 4 are thus qualitative findings about \emph{reported construction concerns}, not a measured distribution of rework events.

\section{Cost and Time Attribution}
Where a verified case has supporting measurements, the additional direct cost may be represented as
\begin{equation}
C_{e}=\sum_{i=1}^{n_e}Q_{ie}R_{ie}+L_{e}+E_{e}+D_{e}+O_{e},
\end{equation}
where $Q_{ie}R_{ie}$ represents verified incremental material or measured work, $L_e$ denotes incremental labour not already included in measured work rates, $E_e$ additional equipment, $D_e$ removal/disposal and $O_e$ other separately supported direct expenditures. Overlapping components must not be double-counted. This is an accounting structure, not a calculated result from the supplied narrative.

Schedule consequences require dated progress and programme records distinguishing corrective work duration from actual project-completion delay. Additional hours at an isolated element do not necessarily extend the controlling path of the project. If resource and programme records are unavailable, the relevant impact is classified \emph{not quantified}, not zero.

\section{Credibility, Ethics and Limitations}
The identity and professional experience of the account's author require verification before student attribution or formal acknowledgement. Site and client names, project addresses and identifiable photographs should be withheld from publication pending appropriate permission and anonymisation. Alternative physical explanations should be considered rather than presenting the first plausible mechanism as a demonstrated cause. The absence of authenticated project case files, contemporaneous site notes, photographic permissions, measured crack data and repair documentation prevents calculation of local case frequencies, causal rankings, financial totals or schedule delays. These limits are observed while still permitting substantive technical interpretation of the supplied concerns.
